В~ходе выполнения данной курсовой работы были успешно решены все поставленные задачи и~достигнута ключевая цель исследования~--- систематизация теоретических сведений об~алгоритмах на графах и~их практическая реализация.

В~теоретической части работы были детально рассмотрены математические основы представления ориентированных графов в~виде матриц смежности, а~также проанализированы ключевые подходы к~поиску кратчайших путей и обходу графовых структур. Был проведён подробный сравнительный анализ алгоритмов обхода в~ширину (BFS), Дейкстры и~Беллмана-Форда по~временной и~пространственной сложности. Выявлены их границы применимости и~ограничения, связанные с~наличием отрицательных весов рёбер.

В~практической части работы была разработана эффективная и~чистая программная реализация алгоритмов на~языке программирования C++. Тестирование разработанного программного комплекса подтвердило его корректность и~высокую скорость работы на различных входных данных. Полученные результаты могут быть использованы при проектировании и~оптимизации реальных информационных систем, использующих графовые модели для маршрутизации и логистики.
